% B2: esquema qualitativo, sem proporções ou comparação de desempenho.
\begin{tikzpicture}[
    caixa/.style={draw=black!65, rounded corners=2pt, align=left,
        text width=4.0cm, minimum height=3.7cm, inner sep=8pt,
        font=\small}, node distance=5mm]
\node[caixa,fill=blue!4] (intrinseca)
    {\textbf{Carga intrínseca}\\[5pt]
    Complexidade dos conceitos e das relações entre os elementos do conteúdo.};
\node[caixa,fill=red!4,right=of intrinseca] (extrinseca)
    {\textbf{Carga extrínseca}\\[5pt]
    Esforço associado à apresentação do material e às atividades que desviam a atenção do conteúdo.};
\node[caixa,fill=green!4,right=of extrinseca] (aprendizagem)
    {\textbf{Esforço de aprendizagem}\\[5pt]
    Recursos dedicados à construção e à organização de esquemas, na formulação clássica de 1998.};
\end{tikzpicture}
